\chapter[METODOLOGIA]{Metodologia}
\label{cap:metodologia}

Este capítulo descreve como o sistema foi construído e avaliado. São
apresentados a caracterização da pesquisa, o procedimento de coleta, as
heurísticas de classificação na sua versão original, os conjuntos de dados, as
métricas e o procedimento experimental de calibração. A arquitetura do sistema
e a sequência de etapas do \textit{pipeline} são descritas no
Capítulo~\ref{cap:projeto}. Os resultados obtidos com esse procedimento são
apresentados no Capítulo~\ref{cap:resultados}, e as conclusões, limitações e
trabalhos futuros, no Capítulo~\ref{cap:conclusao}.

% -----------------------------------------------------------------------------
\section{Caracterização da pesquisa}
\label{sec:caracterizacao}

Quanto à natureza, esta é uma pesquisa aplicada, pois seu objetivo é produzir e
avaliar um sistema que identifica, na Web aberta, editais de fomento com
inscrições em aberto \cite{gil2017}. A abordagem é quantitativa: o desempenho do
sistema é medido por métricas de classificação calculadas sobre amostras
rotuladas manualmente.

Quanto ao delineamento, o estudo é experimental, baseado em \textit{Design Science Research} (DSR), e segue o ciclo de construção e
avaliação de artefatos computacionais descrito por
\citeonline{hevner2004}. Primeiro, implementou-se um \textit{pipeline} baseado em
heurísticas. Em seguida, os parâmetros dessas heurísticas foram calibrados
sobre um conjunto de calibração. Por fim, a configuração resultante foi
congelada e avaliada uma única vez sobre um conjunto de teste independente,
coletado em outra data e mantido fora de todas as decisões de calibração.

A separação entre calibração e teste é o principal controle metodológico do
trabalho. Como as amostras são pequenas (Seção~\ref{sec:dados}), ajustar
parâmetros e medir o desempenho nos mesmos dados superestimaria a qualidade do
sistema \cite{cawley2010}. Por isso, toda escolha de parâmetro, regra ou
variante foi feita exclusivamente com o conjunto de calibração.

% -----------------------------------------------------------------------------
\section{Coleta de dados e reprodutibilidade}
\label{sec:coleta}

O \textit{pipeline} do sistema separa a coleta, que consulta a API de busca e
baixa as páginas, do processamento, que aplica a elas a cadeia de filtros
(Seção~\ref{sec:arquitetura}). Cada execução da coleta gera um
\textit{snapshot}: um diretório que contém o
conteúdo bruto de cada página, nomeado pelo \textit{hash} SHA-256 do próprio
conteúdo, e um arquivo de manifesto. O manifesto registra o instante da coleta,
o identificador do \textit{commit} do código-fonte no Git \cite{git}, as versões
das bibliotecas de extração, as consultas executadas, a configuração de
parâmetros vigente e os metadados de cada item (URL, título, resumo, código de
\textit{status} HTTP e tipo de conteúdo).

O \textit{snapshot} torna os experimentos reprodutíveis. Como a Web muda
continuamente, uma página pode ser alterada ou removida entre duas execuções. Ao
reprocessar o \textit{snapshot}, o sistema opera sempre sobre o mesmo conteúdo.
Além disso, a data de referência usada para decidir se um prazo está vencido é
a data registrada no manifesto, e não a data em que o código é executado.
Assim, o reprocessamento de uma coleta antiga reproduz exatamente as decisões
que o sistema tomaria no dia da coleta.

Para garantir que as alterações no código não mudassem o comportamento do
sistema de forma inadvertida, adotou-se um teste de regressão. Após cada
refatoração, o \textit{snapshot} de calibração era reprocessado e o arquivo de
predições resultante era comparado, byte a byte, com o arquivo de referência.
Toda modificação estrutural foi introduzida como uma opção desligada por padrão
(Seção~\ref{subsec:etapaD}), de modo que o teste devesse resultar em
identidade. Uma diferença indicaria efeito colateral não intencional.

Os \textit{snapshots}, os rótulos e os arquivos de resultado foram versionados
junto com o código, de modo que cada resultado pode ser associado ao
\textit{commit} que o produziu.

% -----------------------------------------------------------------------------
\section{Heurísticas de classificação}
\label{sec:heuristicas}

As três decisões centrais do \textit{pipeline} são tomadas por heurísticas de
pontuação: somas ponderadas de indícios textuais, comparadas a um limiar.
Optou-se por heurísticas, e não por um classificador treinado, por dois
motivos. Primeiro, o volume de dados rotulados é insuficiente para treinar um
modelo supervisionado com segurança, e heurísticas simples podem ter
desempenho competitivo em relação a modelos estatísticos nessas condições
\cite{buckmann2017}. Segundo, as regras são interpretáveis, o que permite
justificar cada decisão do sistema.

Todos os pesos e limiares foram centralizados em uma estrutura de configuração
imutável, o que permite variar um parâmetro sem alterar o código. Esta seção
descreve as heurísticas na \textbf{configuração original}, ponto de partida da
calibração. Os valores alterados pela calibração são apresentados na
Seção~\ref{sec:configuracao}.

Nas expressões a seguir, $\mathbf{1}[\cdot]$ denota a função indicadora, que
vale 1 quando a condição é verdadeira e 0 caso contrário.

\subsection{Filtro de listagem}
\label{subsec:listagem}

Uma página é classificada como listagem quando sua pontuação $L$ atinge o
limiar $\theta_L$. A pontuação soma indícios da URL e do conteúdo, conforme a
Equação~\eqref{eq:listagem}:

\begin{equation}
\label{eq:listagem}
\begin{split}
L = {} & w_k \,\mathbf{1}[\text{URL contém palavra de listagem}]
       + w_p \,\mathbf{1}[\text{URL contém padrão de listagem}] \\
     & + w_d \,\mathbf{1}[n_{\mathrm{datas}} > d_{\mathrm{low}}]
       + w_d \,\mathbf{1}[n_{\mathrm{datas}} > d_{\mathrm{high}}]
       + w_l \,\mathbf{1}[n_{\mathrm{det}} < m] \\
     & + w_e \,\mathbf{1}[\text{texto contém ``editais''}]
       + w_c \,\mathbf{1}[\text{texto contém ``chamadas públicas''}]
\end{split}
\end{equation}

As palavras de URL incluem termos como ``editais'', ``noticias'', ``chamadas''
e ``resultados''. Os padrões de URL são segmentos de caminho, como
\texttt{/noticias/}, \texttt{/editais/}, \texttt{/chamadas/} e
\texttt{/resultados/}. $n_{\mathrm{datas}}$ é o número de datas no formato
DD/MM/AAAA presentes no texto: muitas datas sugerem uma lista de chamadas.
$n_{\mathrm{det}}$ é o número de ocorrências de termos típicos de um edital
individual, como ``edital'', ``prazo'', ``inscrição'' e ``cronograma'': poucas
ocorrências sugerem uma página genérica. Os dois últimos termos captam plurais
característicos de páginas agregadoras. A configuração original continha ainda
uma regra baseada na contagem de \textit{links} no texto, com três parâmetros
próprios. Os valores originais estão na Tabela~\ref{tab:listagem_original}.

\begin{table}[htbp]
\centering
\caption{Parâmetros do filtro de listagem na configuração original}
\label{tab:listagem_original}
\small
\begin{tabular}{l l c}
\toprule
Símbolo & Significado & Valor \\
\midrule
$w_k$               & Peso da palavra de listagem na URL      & 2 \\
$w_p$               & Peso do padrão de listagem na URL       & 3 \\
$w_d$               & Peso de cada faixa de contagem de datas & 2 \\
$d_{\mathrm{low}}$  & Primeira faixa de contagem de datas     & 5 \\
$d_{\mathrm{high}}$ & Segunda faixa de contagem de datas      & 10 \\
$w_l$               & Peso da baixa ocorrência de termos de edital & 2 \\
$m$                 & Mínimo de termos de edital              & 2 \\
$w_e$               & Peso do plural ``editais''              & 1 \\
$w_c$               & Peso do plural ``chamadas públicas''    & 1 \\
$\theta_L$          & Limiar de classificação                 & 5 \\
\bottomrule
\end{tabular}
\fonte{Elaborado pelo autor.}
\end{table}

\subsection{Filtro de relevância}
\label{subsec:relevancia}

A pontuação de relevância $R$ conta ocorrências de palavras inteiras de um
vocabulário ponderado, no texto convertido para minúsculas, conforme a
Equação~\eqref{eq:relevancia}:

\begin{equation}
\label{eq:relevancia}
R = \sum_{k \in K^{+}} w_k\, n_k + \sum_{k \in K^{-}} w_k\, n_k + B(\rho),
\qquad
\rho = \frac{\sum_{k \in K^{+}} n_k}{N}
\end{equation}

O vocabulário positivo $K^{+}$ contém ``edital'' e ``chamada pública''
(peso 3) e ``financiamento'', ``inovação'' e ``incentivo'' (peso 2). O
vocabulário negativo $K^{-}$ contém ``resultado'' e ``notícia'' (peso $-2$),
além de ``publicado'' e ``evento'' (peso $-1$). $n_k$ é o número de
ocorrências do termo $k$ e $N$ é o número total de palavras. O bônus de
densidade é definido pela Equação~\eqref{eq:bonus}, e a página é considerada
relevante quando $R \geq 3$. A configuração original tinha ainda um limiar
secundário, de valor 2, aplicado após o limiar principal.

\begin{equation}
\label{eq:bonus}
B(\rho) =
\begin{cases}
3, & \rho > 0{,}05 \\
1, & 0{,}02 < \rho \leq 0{,}05 \\
0, & \text{caso contrário}
\end{cases}
\end{equation}

\subsection{Extração de prazo}
\label{subsec:prazo}

Na configuração original, a extração de prazo opera em quatro passos.

\begin{enumerate}
  \item \textbf{Frases candidatas}: o texto é segmentado em frases nos pontos
    finais e quebras de linha, e são mantidas as frases que contêm ao menos um
    dos 18 termos indicativos de prazo, como ``prazo'', ``inscrições'',
    ``submissão'', ``até'', ``data limite'', ``encerramento'' e
    ``cronograma''.
  \item \textbf{Pontuação da frase}: cada frase candidata recebe a pontuação
    $S$ da Equação~\eqref{eq:frase}, e são descartadas as frases com
    $S \leq 0$.
    \begin{equation}
    \label{eq:frase}
    S = 2\, n_{\mathrm{pos}} - 3\, n_{\mathrm{neg}}
        + 3\,\mathbf{1}[\text{frase contém ``até''}]
        + 4\,\mathbf{1}[\text{frase contém ``encerramento''}]
    \end{equation}
    Nela, $n_{\mathrm{pos}}$ é o número de termos indicativos distintos presentes
    na frase e $n_{\mathrm{neg}}$ é o número de termos negativos distintos, como
    ``resultado'', ``divulgação'', ``homologação'' e ``assinatura'', que
    indicam etapas posteriores à inscrição.
  \item \textbf{Data da frase}: a frase inteira é entregue à biblioteca
    \textit{dateparser} \cite{dateparser}, que tenta interpretá-la como uma data.
  \item \textbf{Decisão}: as datas anteriores à data da coleta são separadas
    como vencidas. Entre as datas restantes, a de maior pontuação é o prazo
    estimado e o item é aceito. Se só houver datas vencidas, o edital é
    considerado encerrado e o item é descartado.
\end{enumerate}

Quando nenhuma frase candidata fornece uma data futura, o sistema recorre a um
procedimento alternativo (\textit{fallback}): busca no texto completo datas nos
formatos numéricos, por expressão regular, interpreta-as com a biblioteca
\textit{dateparser} e atribui a cada uma a pontuação 1. Persistindo a ausência de prazo,
a extração é repetida sobre os primeiros 2.000 caracteres do texto.

% -----------------------------------------------------------------------------
\section{Conjuntos de dados e protocolo de rotulação}
\label{sec:dados}

Foram construídos dois conjuntos rotulados, cada um a partir de uma coleta
distinta (Tabela~\ref{tab:conjuntos}). O conjunto de calibração foi usado em
todas as decisões de ajuste. O conjunto de teste foi reservado e consultado uma
única vez, após o congelamento da configuração. Antes disso, admitiu-se apenas
a verificação de formato e de consistência dos seus rótulos, sem cálculo de
métricas.

\begin{table}[htbp]
\centering
\caption{Composição dos conjuntos rotulados}
\label{tab:conjuntos}
\small
\setlength{\tabcolsep}{4pt}
\begin{tabular}{l c c c c c c c c}
\toprule
Conjunto & Coleta & \makecell{Itens\\coletados} & \makecell{Itens\\rotulados}
  & Listagens & Editais & \makecell{Editais\\com data} & \makecell{Editais\\abertos}
  & Relevantes \\
\midrule
Calibração & 27/10/2026 & 47 & 47 & 18 & 20 & 19 & 1 & 27 \\
Teste      & 07/11/2026 & 45 & 26 &  5 & 14 & 12 & 3 & 10 \\
\bottomrule
\end{tabular}
\fonte{Elaborado pelo autor.}
\end{table}

Dos 45 itens obtidos na segunda coleta, 19 já constavam da coleta de
calibração e foram descartados, para que nenhum item usado na calibração fosse
reutilizado na avaliação. Os 26 itens restantes foram todos rotulados, sem
amostragem adicional. Como os itens repetidos eram, em geral, páginas
institucionais estáveis, como as listagens de chamadas das agências de
fomento, o conjunto de teste tem proporcionalmente menos listagens que o de
calibração.

A rotulação foi manual e seguiu um guia escrito, aplicado aos dois conjuntos
para que fossem rotulados com os mesmos critérios. Cada item recebeu os rótulos
descritos na Tabela~\ref{tab:rotulos}.

\begin{table}[htbp]
\centering
\caption{Rótulos atribuídos a cada item}
\label{tab:rotulos}
\small
\begin{tabularx}{\textwidth}{>{\ttfamily}l X l}
\toprule
\normalfont Rótulo & Pergunta respondida & Valores \\
\midrule
is\_listing\_real       & A página lista vários editais ou chamadas? & sim / não \\
is\_edital\_real        & A página é um edital ou a página oficial de um edital? & sim / não \\
is\_relevant\_real      & A página é relevante para o escopo do sistema (tecnologia e inovação)? & sim / não \\
deadline\_real          & Existe prazo de inscrição explícito na página? & sim / não \\
deadline\_status\_real  & A inscrição estava aberta na data da coleta? & sim / não \\
deadline\_date\_real    & Qual a data máxima de inscrição? & \makecell[l]{AAAA-MM-DD,\\ou vazio} \\
\bottomrule
\end{tabularx}
\fonte{Elaborado pelo autor.}
\end{table}

O guia estabelece duas regras de consistência. Primeira: os rótulos
\texttt{is\_listing\_real} e \texttt{is\_edital\_real} são mutuamente
exclusivos, pois uma página que lista editais não é, ela própria, um edital.
Segunda: a combinação \texttt{deadline\_real}~=~não e
\texttt{deadline\_status\_real}~=~sim identifica uma chamada de fluxo contínuo,
aberta sem data-limite explícita. Os rótulos de prazo só têm significado para
editais; para os demais itens, são tratados como negativos na avaliação
(Seção~\ref{subsec:derivados}).

O conjunto de calibração tem um único edital aberto, e o de teste tem três.
As métricas do resultado final do sistema, que dependem desses itens, têm por
isso variância muito alta. Por esse motivo, as decisões de calibração
apoiaram-se em métricas por etapa, calculadas sobre os 47 itens ou sobre os 19
editais com data.

% -----------------------------------------------------------------------------
\section{Métricas de avaliação}
\label{sec:metricas}

\subsection{Métricas de classificação}
\label{subsec:classificacao}

Cada decisão binária do sistema foi comparada ao rótulo manual correspondente,
produzindo uma matriz de confusão com verdadeiros positivos ($VP$), falsos
positivos ($FP$), falsos negativos ($FN$) e verdadeiros negativos ($VN$). A
partir dela foram calculadas a precisão $P$ e a revocação $R$, definidas no
Referencial Teórico, e a medida $F_\beta$, que as combina em um único valor
\cite{vanrijsbergen1979,manning2008,sokolova2009}:

\begin{equation}
\label{eq:fbeta}
P = \frac{VP}{VP + FP},
\qquad
R = \frac{VP}{VP + FN},
\qquad
F_{\beta} = \frac{(1 + \beta^{2})\, P\, R}{\beta^{2}\, P + R}
\end{equation}

O parâmetro $\beta$ pondera a revocação em relação à precisão. Foram reportados
$F_{0{,}5}$, $F_1$ e $F_2$, e a medida-alvo de cada etapa foi escolhida conforme
o custo de seus erros na cadeia:

\begin{itemize}
  \item \textbf{Filtro de listagem}: $F_{0{,}5}$. Um falso positivo dessa etapa é
    um edital classificado como listagem e, portanto, descartado sem recurso. A
    precisão é priorizada.
  \item \textbf{Filtro de relevância}: $F_2$. Um item irrelevante aprovado
    ainda pode ser barrado pela extração de prazo, enquanto um item relevante
    rejeitado é perdido. A revocação é priorizada.
  \item \textbf{Extração de prazo}: acurácia da data, definida como a fração
    dos editais com data rotulada cuja data estimada coincide exatamente com o
    rótulo \texttt{deadline\_date\_real}.
\end{itemize}

\subsection{Métricas operacionais}
\label{subsec:operacionais}

As medidas acima avaliam cada etapa isoladamente. Para avaliar o efeito de uma
mudança sobre o sistema completo, definiram-se três contagens complementares:

\begin{itemize}
  \item \textbf{Editais perdidos}: número de itens rotulados como edital que o
    filtro de listagem descarta.
  \item \textbf{Itens errados aceitos}: número de itens aceitos pelo sistema
    completo cujo rótulo de resultado final é negativo, ou seja, os falsos
    positivos do resultado final.
  \item \textbf{Aceitos com a data certa}: entre os itens aceitos, quantos têm
    o prazo estimado igual ao rotulado.
\end{itemize}

\subsection{Rótulos derivados}
\label{subsec:derivados}

Dois rótulos foram derivados dos rótulos manuais. O rótulo de prazo aberto de
um item é negativo quando o item não é edital. Para os editais, é o valor de
\texttt{deadline\_status\_real} quando informado, e é negativo quando o edital
não tem prazo explícito. O rótulo de resultado final é positivo somente quando
o item é edital, é relevante e tem inscrições abertas na data da coleta, que é
exatamente o que o sistema deve entregar. Itens sem informação suficiente para
derivar o rótulo foram excluídos da métrica correspondente.

\subsection{Modos de avaliação}
\label{subsec:modos}

A avaliação foi feita de duas formas. Na avaliação por etapa, todo item passa
por todas as heurísticas, independentemente da cadeia, o que permite medir cada
etapa sobre a amostra inteira. Na avaliação em cadeia, reproduz-se a sequência
de filtros do \textit{pipeline}, e uma etapa não executada para um item conta
como decisão negativa. A primeira forma orienta a calibração de cada
heurística. A segunda mede o que o usuário do sistema efetivamente recebe.

% -----------------------------------------------------------------------------
\section{Procedimento de calibração}
\label{sec:calibracao}

A calibração foi organizada em cinco etapas sequenciais, de A a E. Cada etapa
só começou após a conclusão da anterior, e o conjunto de teste só foi usado na
última.

\subsection{Etapa A: correção da avaliação}
\label{subsec:etapaA}

Antes de ajustar qualquer parâmetro, revisou-se o próprio procedimento de
avaliação, pois uma métrica mal definida conduziria a calibração na direção
errada. Nessa revisão, definiram-se os rótulos derivados da
Seção~\ref{subsec:derivados}, de modo que itens que não são editais passassem a
contar como negativos nas métricas de prazo e de resultado final, em vez de
serem excluídos. Também foram introduzidas as métricas operacionais da
Seção~\ref{subsec:operacionais} e a medida $F_{0{,}5}$.

\subsection{Etapa B: conjunto de teste}
\label{subsec:etapaB}

Realizou-se a segunda coleta, descartaram-se os itens já presentes na coleta
de calibração e rotularam-se os restantes com o mesmo guia
(Seção~\ref{sec:dados}). A partir desse momento, o conjunto foi tratado como
reservado: nenhuma métrica foi calculada sobre ele até a Etapa E.

\subsection{Etapa C: análise de sensibilidade}
\label{subsec:etapaC}

Aplicou-se uma análise de sensibilidade do tipo um-fator-por-vez
(\textit{one-at-a-time}, OAT) \cite{saltelli2008}. Cada um dos 36 parâmetros
numéricos foi variado ao longo de uma grade de valores, mantendo-se os demais
na configuração de base, e registrou-se a métrica-alvo do grupo do parâmetro
(Seção~\ref{subsec:classificacao}) para cada valor. Os parâmetros abrangeram
pesos e limiares dos três filtros, inclusive o peso individual de cada termo do
vocabulário de relevância. Para cada parâmetro foram calculados:

\begin{itemize}
  \item a \textbf{amplitude}, diferença entre o maior e o menor valor da
    métrica-alvo ao longo da grade, que mede a influência do parâmetro;
  \item a amplitude da medida $F_1$ do resultado final, que mede a influência
    sobre o sistema completo;
  \item o \textbf{valor sugerido}, definido como o valor central da faixa de
    valores empatados no melhor resultado. A escolha do centro do patamar, e não
    de sua extremidade, privilegia a robustez, pois o resultado não depende de
    um valor isolado.
\end{itemize}

Para os filtros de listagem e de prazo, o melhor resultado foi definido por um
critério lexicográfico, e não apenas pela métrica-alvo. Uma configuração é
preferível a outra quando, nesta ordem: (1) não aumenta o número de itens
errados aceitos em relação à base; (2) perde menos editais; (3) tem maior valor
da métrica-alvo. O critério reflete o objetivo do sistema: uma configuração que
melhora a $F_{0{,}5}$ do filtro de listagem à custa de descartar editais não é
desejável.

Além de orientar o ajuste, a análise teve função diagnóstica. Todo parâmetro de
amplitude nula foi investigado no código-fonte, para distinguir parâmetros que
apenas não têm influência na amostra de defeitos de implementação que impedem
o parâmetro de atuar.

\subsection{Etapa D: variantes estruturais, busca em grade e validação}
\label{subsec:etapaD}

As correções e hipóteses de melhoria levantadas na Etapa C foram implementadas
como variantes estruturais. Cada variante é uma opção de configuração desligada
por padrão, o que mantém o comportamento original intacto (verificado pelo
teste de regressão) e permite comparar a variante com a base sobre o mesmo
\textit{snapshot}. Uma variante foi aceita somente quando atendia ao critério
lexicográfico da Etapa C. Mecanismos sem efeito em nenhum valor testado foram
removidos, desde que a remoção mantivesse as métricas idênticas.

Incorporadas as variantes aceitas, a análise OAT foi repetida sobre a nova
base. Como a análise OAT não captura interações entre parâmetros, realizou-se
em seguida uma busca em grade restrita ao filtro de listagem, com os seis
parâmetros e valores da Tabela~\ref{tab:grade}, em um total de 432
combinações. As combinações foram ordenadas pelo critério lexicográfico. Entre
combinações empatadas, adotou-se a que altera menos parâmetros em relação à
base, segundo o princípio da parcimônia.

\begin{table}[htbp]
\centering
\caption{Valores avaliados na busca em grade do filtro de listagem}
\label{tab:grade}
\small
\begin{tabular}{l l}
\toprule
Parâmetro & Valores \\
\midrule
$w_p$               & 0, 1, 3 \\
$w_k$               & 2, 3 \\
$w_d$               & 2, 3, 4 \\
$d_{\mathrm{high}}$ & 10, 20, 30 \\
$w_e$               & 1, 2 \\
$\theta_L$          & 3, 4, 5, 6 \\
\bottomrule
\end{tabular}
\fonte{Elaborado pelo autor.}
\end{table}

Como 432 combinações avaliadas sobre 47 itens podem produzir um ótimo devido ao
acaso \cite{cawley2010}, a escolha da grade foi submetida a duas verificações:

\begin{itemize}
  \item \textbf{Validação cruzada} \cite{kohavi1995} com 5 partições e 10
    repetições, totalizando 50 partições. Em cada uma, a busca em grade foi
    refeita usando apenas os rótulos de treino, e a configuração escolhida foi
    avaliada nos itens retidos. Compararam-se a configuração escolhida e a base
    pela soma, sobre as partições, dos editais perdidos e dos itens errados
    aceitos. A estabilidade da escolha foi medida pela frequência com que a
    mesma combinação era selecionada.
  \item \textbf{\textit{Bootstrap}} \cite{efron1993} com 2.000 reamostragens com
    reposição dos itens, para estimar o intervalo de confiança de 95\% da
    diferença de editais perdidos e de itens errados aceitos entre a
    configuração calibrada e a original. A reamostragem é feita sobre os itens,
    e não sobre as métricas, de modo que cada reamostra compara as duas
    configurações sobre os mesmos itens.
\end{itemize}

\subsection{Etapa E: congelamento e teste}
\label{subsec:etapaE}

Ao fim da Etapa D, a configuração calibrada tornou-se a configuração padrão do
sistema e foi registrada em um \textit{commit} do repositório. O
\textit{snapshot} de calibração foi reprocessado para gerar o novo arquivo de
referência do teste de regressão. Só então o conjunto de teste foi avaliado,
uma única vez, com a configuração original e com a calibrada. Como ambas foram
fixadas antes de qualquer consulta ao conjunto de teste, a comparação entre
elas é legítima. Nenhum ajuste foi feito após essa avaliação, e o
\textit{bootstrap} da Etapa D foi repetido sobre o conjunto de teste. Os
comandos para repetir as cinco etapas sobre coletas novas estão descritos na
Seção~\ref{sec:impl_recalibracao}.
